\section{Nonuniqueness with a zero set of empty interior}
\label{sec:numa-thin-zero-set}

We now keep the Hermitian background of \cref{thm:degenerate-hermitian-nonuniqueness}, but change the density so that it is positive off $M_-\cup M_+$, where $M_\pm=\{r=\pm1\}$. For potentials of the form $G(r)-\rho$, the equation on the middle band becomes a first-order ordinary differential equation for the slope $G'$. Barriers keep its solutions between the admissible slopes, and infinite-order vanishing at the endpoints then makes their extensions by zero smooth.

\Needspace{9\baselineskip}
\begin{theorem}\label{thm:numa-thin-zero-set}\index{Zero set!thin degeneracy set}
There are a smooth nonclosed Hermitian form $\omega$ on $\mathbb P^3$ and a smooth nonnegative density $g$, positive off two disjoint smooth real hypersurfaces, such that the equation
\[
 \mathrm{MA}_\omega(u)=g\omega_X^3
\]
has uncountably many distinct smooth solutions normalized by $\sup_Xu=0$. The background $\omega$ has bounded mass.
\end{theorem}

Thus a zero set with empty interior, even one of measure zero, does not ensure uniqueness of smooth normalized solutions.

\Needspace{12\baselineskip}
We use the coefficients of \cref{subsec:numa-scalar-coefficients}, with the explicit choice
\[
 p(r)=C\exp\!\left(-\frac1{1-r^2}\right),\qquad -1<r<1,
 \qquad \int_{-1}^1p(r)\,\mathrm{d}r=1,
\]
extended by zero, where the integral condition determines $C>0$; then $q=p'$. Fix $0<s<1$, and set
\[
 k_s(r)=(x(r)+sp(r))(y(r)-sp(r))(z(r)+sq(r)).
\]
For $|r|>1$ this is $xyz$. If $G'=sp$, then the potential $G(r)-\rho$ has $\omega+\ddc(G(r)-\rho)=\omega_0+sH$. This form is strictly positive for $-1<r<1$, since
\[
 x+sp=a+(1+s)p>0,\quad
 y-sp=b+(1-s)p>0,\quad
 z_0+sq\geq(1-s)z_0>0.
\]
It is also strictly positive outside the closed band, and at $r=\pm1$ it equals the rank-one form $\omega_0$. Consequently
\begin{equation}\label{eq:numa-thin-density}
 g\omega_X^3=(\omega_0+sH)^3
\end{equation}
defines a smooth density $g$ that is positive exactly off $M_-\cup M_+$. These hypersurfaces have measure zero and empty interior. On the band $|r|<1$, a potential $G(r)-\rho$ with slope $h=G'$ solves the equation exactly when $(x+h)(y-h)(z+h')=k_s$, which is the differential equation \eqref{eq:numa-thin-ode} below.

\Needspace{14\baselineskip}
\begin{lemma}\label{lem:numa-thin-trapping}
There exists $\eta_0>0$ such that for every $0<\eta<\eta_0$ the initial-value problem
\begin{equation}\label{eq:numa-thin-ode}
 h'=\frac{k_s(r)}{(x(r)+h)(y(r)-h)}-z(r),
 \qquad h(0)=sp(0)-\eta,
\end{equation}
has a smooth solution on $(-1,1)$ with $-x<h<y$ and $h<sp$. Moreover, on some fixed endpoint neighborhoods there are constants $0<c_L<s<t<1$, independent of $\eta$, such that
\[
 c_Lp<h<sp\quad\hbox{near }-1,
 \qquad -tp<h<sp\quad\hbox{near }1.
\]
\end{lemma}
\begin{proof}
The function $h_s=sp$ solves \eqref{eq:numa-thin-ode} and lies strictly between $-x$ and $y$ throughout $(-1,1)$. Write $F_s(r,h)$ for the right-hand side of the equation, which is smooth on the open domain $-x(r)<h<y(r)$.

Near $-1$ the coefficients are exactly
\[
 x=1+p,\qquad y=p,\qquad z=q>0.
\]
For a constant $c\in(-1,1)$, the product of coefficients for the trial slope $cp$ is
\[
 k_c=(1-c^2)\bigl(1+(1+c)p\bigr)pq.
\]
Fix $c_L\in(0,s)$. At $p=0$, the factor preceding $pq$ for $c=c_L$ exceeds that for $c=s$ by $s^2-c_L^2>0$. Thus, on a sufficiently small fixed neighborhood of the left endpoint, $k_{c_L}>k_s$ and
\begin{equation}\label{eq:numa-left-trap}
 F_s(r,c_Lp)-(c_Lp)'=
 \frac{k_s-k_{c_L}}{(x+c_Lp)(y-c_Lp)}<0.
\end{equation}
Thus $c_Lp$ is a lower barrier when the equation is solved backward, that is, toward $-1$.

Near $1$ the exact coefficients are
\[
 x=p,\qquad y=1+p,\qquad z=-q>0,
\]
and the trial slope $cp$ has product of coefficients
\[
 k_c=(1-c^2)\bigl(1+(1-c)p\bigr)p(-q).
\]
Fix $s<t<1$. At $p=0$, the factor for $c=s$ exceeds that for $c=-t$ by $t^2-s^2>0$. On a sufficiently small fixed neighborhood of the right endpoint, therefore, $k_s>k_{-t}$, so
\begin{equation}\label{eq:numa-right-trap}
 F_s(r,-tp)-(-tp)'=
 \frac{k_s-k_{-t}}{(x-tp)(y+tp)}>0.
\end{equation}
Thus $-tp$ is a lower barrier when the equation is solved forward, toward $1$.

Choose $r_L<0<r_R$ in these respective neighborhoods. The graph of $h_s$ over $[r_L,r_R]$ is a compact subset of the domain of $F_s$. Existence and continuous dependence on initial values for smooth ordinary differential equations therefore give $\eta_0>0$ such that, for every $0<\eta<\eta_0$, the solution with initial value $sp(0)-\eta$ exists on $[r_L,r_R]$, stays in the domain, and satisfies
\[
 h(r_L)>c_Lp(r_L),\qquad h(r_R)>-tp(r_R).
\]
By local uniqueness, two solutions of the equation cannot cross. Since $h(0)<h_s(0)$, we have $h<h_s$ wherever $h$ exists.

Follow the solution backwards from $r_L$; it cannot cross $c_Lp$. Indeed, at a first crossing point the difference $h-c_Lp$ would vanish and be positive immediately to its right, so its derivative there would be nonnegative, whereas \eqref{eq:numa-left-trap} shows that this derivative is strictly negative. Thus $c_Lp<h<sp$ on the part of the interval of existence to the left of $r_L$. Likewise, \eqref{eq:numa-right-trap} prevents a first downward crossing of $-tp$ to the right of $r_R$, giving $-tp<h<sp$ there.

These bounds also show that the maximal interval of existence is all of $(-1,1)$. Indeed, near a possible endpoint inside $(-1,1)$, the factors $x+h$ and $y-h$ stay bounded away from zero, so the graph of $h$ stays in a compact subset of the domain of $F_s$, and the continuation theorem for ordinary differential equations extends the solution beyond that endpoint. Hence the solution exists on $(-1,1)$ and satisfies the claimed bounds.
\end{proof}

\Needspace{7\baselineskip}
\begin{lemma}\label{lem:numa-thin-flat-extension}
Every solution furnished by \cref{lem:numa-thin-trapping}, extended by zero outside $(-1,1)$, is smooth on $\mathbb R$ and vanishes to infinite order at $\pm1$.
\end{lemma}
\begin{proof}
Let $h$ be one of the solutions in \cref{lem:numa-thin-trapping}. On either endpoint neighborhood put $\xi=h/p$; since $p>0$ on $(-1,1)$, this quotient is smooth there, and the trapping bounds place $\xi$ in a fixed compact subinterval of $(-1,1)$. Define the rational functions
\[
 R_L(p,\xi)=
 \frac{(1-s^2)(1+(1+s)p)}{(1+(1+\xi)p)(1-\xi)}-1,
\]
\[
 R_R(p,\xi)=1-
 \frac{(1-s^2)(1+(1-s)p)}{(1+\xi)(1+(1-\xi)p)}.
\]
The explicit formulas for the coefficients near the endpoints give
\[
 h'=q\,R_L(p,\xi)\quad\hbox{on the left},\qquad
 h'=q\,R_R(p,\xi)\quad\hbox{on the right}.
\]
Both rational functions are smooth, with all derivatives bounded, on a neighborhood of the relevant set of pairs $(p,\xi)$ with $p\geq0$ small and $\xi$ in the compact interval above. Writing $R_{L/R}$ for the function corresponding to the endpoint under consideration, we obtain from $h'=p'\xi+p\xi'$ and $q=p'$
\begin{equation}\label{eq:numa-scaled-ode}
 \xi'=\frac{p'}p\bigl(R_{L/R}(p,\xi)-\xi\bigr).
\end{equation}
For the chosen bump,
\[
 \frac{p'}p=-\frac{2r}{(1-r^2)^2}.
\]
This rational function and each of its derivatives grow at most like a power of the inverse distance to the endpoint. Moreover, each derivative of $p$ is $p$ times a rational function with the same property.

An induction using \eqref{eq:numa-scaled-ode} shows that every derivative $\xi^{(j)}$ grows at most polynomially in the inverse distance to the endpoint. The case $j=0$ is the trapping bound on $\xi$. For the induction step, differentiate the right-hand side of \eqref{eq:numa-scaled-ode}, including its term $-\xi$, $j$ times. Every resulting term is a product of a derivative of $p'/p$, a bounded partial derivative of $R_{L/R}$, and derivatives of $p$ and of $\xi$ of order at most $j$. All of these factors have polynomial bounds, by the induction hypothesis for the derivatives of $\xi$ and since every derivative of $p$ is bounded. This gives the bound for $\xi^{(j+1)}$.

Leibniz' rule for $h=p\xi$ now gives, for each $j$, constants $C_j,N_j$ such that
\[
 |h^{(j)}(r)|\leq C_j p(r)(1-|r|)^{-N_j}
\]
near the endpoints. The right-hand side tends to zero at the endpoints, because $p$ decays exponentially in the inverse distance to the endpoint. Thus all derivatives of $h$ tend to zero at $\pm1$. Extend each derivative by zero outside $(-1,1)$. By the mean value theorem, the difference quotient of $h^{(j)}$ at an endpoint tends to the limit of $h^{(j+1)}$, which is zero. Hence these extensions are the successive derivatives of the extension of $h$, which is therefore smooth and vanishes to infinite order at $\pm1$.
\end{proof}

\begin{proof}[of \cref{thm:numa-thin-zero-set}]
Let $\eta_0$ be as in \cref{lem:numa-thin-trapping}, and for $0<\eta<\eta_0$ denote by $h_\eta$ the solution with initial value $sp(0)-\eta$. Extend $h_\eta$ by zero, which gives a smooth function by the preceding lemma, and take a primitive $G_\eta$. It is constant near each end of the real line, so $G_\eta(r)$ extends smoothly to all of $\mathbb P^3$, including the two omitted projective lines. Set
\[
 \widetilde u_\eta=G_\eta(r)-\rho.
\]
On the open band,
\[
 \omega+\ddc\widetilde u_\eta
 =(x+h_\eta)P+(y-h_\eta)Q+(z+h_\eta')R.
\]
The coefficients of $P$ and $Q$ are strictly positive by \cref{lem:numa-thin-trapping}, while the equation gives
\[
 z+h_\eta'=\frac{k_s}{(x+h_\eta)(y-h_\eta)}>0.
\]
Outside the band the form is $\omega_0$, and on the two boundary hypersurfaces it is also the rank-one form $\omega_0$, since $h_\eta$ vanishes to infinite order there. It is therefore semipositive on $X$, and by \eqref{eq:numa-cube} its cube equals $6k_sP\wedge Q\wedge R=(\omega_0+sH)^3$ on the band and $\omega_0^3=(\omega_0+sH)^3$ elsewhere. Thus its cube is everywhere the measure $g\omega_X^3$ of \eqref{eq:numa-thin-density}. Finally, normalize by subtracting $\sup_X\widetilde u_\eta$.

Distinct values of $\eta$ give different slopes at $r=0$, so the corresponding potentials do not differ by a constant, and their normalized versions are distinct. This gives an uncountable family. The background is that of \cref{thm:degenerate-hermitian-nonuniqueness}, so it has bounded mass.
\end{proof}

\Needspace{8\baselineskip}
The density vanishes to infinite order on the two hypersurfaces. Their complement has three components, and on each of the two outer components, two different solutions differ by a constant, with different constants on the two components. The differential equation on the middle component joins these constants smoothly across the degenerate hypersurfaces. This is why ellipticity at every point of the positive-density set does not give a maximum principle argument across its boundary.
